\documentclass[output=paper,colorlinks,citecolor=brown,arabicfont,georgian]{langscibook}
\ChapterDOI{10.5281/zenodo.18845610}
\author{Zakharia Pourtskhvanidze\orcid{}\affiliation{Goethe-University Frankfurt am Main} and Marina Beridze\orcid{}\affiliation{Javakhishvili State University, Georgia}}
\title{Script and identity: Alphabets, amulets, and anchors}
\abstract{%

The paper, ``Script and identity: Alphabets, amulets, and anchors'', examines 
%\todo[inline]{This paper examines}
the central role of the Georgian alphabet in constructing and preserving identity within isolated Georgian linguistic enclaves, focusing on the Fereydani Georgian community in Iran. The study argues that the Georgian script transcends its function as a linguistic medium and serves as a cultural and religious symbol that anchors collective identity in the midst of geographical and cultural isolation. Separated from the Georgian heartland during the Ottoman-Persian conflicts, the Fereydani Georgians experienced centuries of isolation that fostered the independent development of their language and culture. The Georgian script emerged as a counterpoint to the dominant Arabic script, symbolizing resistance to assimilation and the preservation of ethnic and spiritual heritage. Beyond textual communication, religious texts, amulets and artefacts bearing the Georgian script serve as tangible links to the ancestral homeland, reinforcing cultural memory and solidarity. The article places this phenomenon in a broader historical context, showing how forced relocations fragmented Georgian communities, creating linguistic and cultural islands. Women played a crucial role in preserving written traditions, particularly through the clandestine Dedabruli script developed under Ottoman rule to covertly maintain Orthodox beliefs and cultural practices. Ultimately, the paper argues that the Georgian alphabet is central to the cultural identity of these isolated communities. It transcends practical communication to become a ritual and sacred element that embodies and perpetuates their identity. The Georgian script, the study shows, is not just a linguistic tool, but a lived experience that is integral to the preservation of collective memory and cultural uniqueness.
}


\IfFileExists{../localcommands.tex}{
   \addbibresource{../localbibliography.bib}
   \usepackage{tabularx,multicol}
\usepackage{url}
\urlstyle{same}
\usepackage{textcomp}
 
\usepackage{langsci-optional}
\usepackage{langsci-lgr}
\usepackage{langsci-gb4e}
\usepackage{langsci-branding}

% \usepackage{fontspec}
\babelprovide{georgian} % Georgian language
\babelprovide{armenian} % Armenian language
\babelfont[georgian,armenian]{rm}{DejaVu Sans}
\newfontfamily\georgianfont{DejaVu Sans}
\newfontfamily\armenianfont{DejaVu Sans}

\usepackage{dialogue}
\usepackage{attrib}
\usepackage{attrib} % speaker labels in Koca's chapter
\usepackage{phoenician} % alf, bet in Mengozzi's chapter

   

\makeatletter
\let\thetitle\@title
\let\theauthor\@author
\makeatother

\newcommand{\togglepaper}[1][0]{
   \bibliography{../localbibliography}
   \papernote{\scriptsize\normalfont
     \theauthor.
     \titleTemp.
In Saloumeh Gholami \& Ulrich Mehlem (eds.), \textit{Identity formation, language and migration dynamics: Minorities in the MENA region and the German diaspora}, Berlin: Language Science Press [preliminary page numbering]
   }
   \pagenumbering{roman}
   \setcounter{chapter}{#1}
   \addtocounter{chapter}{-1}
}

\newbool{bookcompile}
\booltrue{bookcompile}
\newcommand{\bookorchapter}[2]{\ifbool{bookcompile}{#1}{#2}}


\newcommand{\georgian}[1]{{\georgianfont #1}}
\newcommand{\armenian}[1]{{\armenianfont #1}}

\newcommand{\textarab}[1]{{\arabicfont \beginR #1} \endR}
\newcommand{\texthebrew}[1]{{\hebrewfont \beginR #1} \endR}
\newcommand{\textsyriac}[1]{{\syriacfont \beginR #1} \endR}


\newfontfamily\phoenicianfont{Tuffy Regular.ttf}

   %% hyphenation points for line breaks
%% Normally, automatic hyphenation in LaTeX is very good
%% If a word is mis-hyphenated, add it to this file
%%
%% add information to TeX file before \begin{document} with:
%% %% hyphenation points for line breaks
%% Normally, automatic hyphenation in LaTeX is very good
%% If a word is mis-hyphenated, add it to this file
%%
%% add information to TeX file before \begin{document} with:
%% %% hyphenation points for line breaks
%% Normally, automatic hyphenation in LaTeX is very good
%% If a word is mis-hyphenated, add it to this file
%%
%% add information to TeX file before \begin{document} with:
%% \include{localhyphenation}
\hyphenation{
    par-a-digm
    Ver-to-vec
    Go-go-lin
    sa-gen
    mo-misch
    ant-wor-te-te
    dix-wa-zim
    spre-chen
    B-Verf-GE
    Swe-dish
    re-mains
    schrei-ben
    di-pey-vin
    Chris-ten-sen
    Wes-sen-dorf
    Ja-va-khi-shvi-li
    ana-lysis
}

\hyphenation{
    par-a-digm
    Ver-to-vec
    Go-go-lin
    sa-gen
    mo-misch
    ant-wor-te-te
    dix-wa-zim
    spre-chen
    B-Verf-GE
    Swe-dish
    re-mains
    schrei-ben
    di-pey-vin
    Chris-ten-sen
    Wes-sen-dorf
    Ja-va-khi-shvi-li
    ana-lysis
}

\hyphenation{
    par-a-digm
    Ver-to-vec
    Go-go-lin
    sa-gen
    mo-misch
    ant-wor-te-te
    dix-wa-zim
    spre-chen
    B-Verf-GE
    Swe-dish
    re-mains
    schrei-ben
    di-pey-vin
    Chris-ten-sen
    Wes-sen-dorf
    Ja-va-khi-shvi-li
    ana-lysis
}

   \boolfalse{bookcompile}
   \togglepaper[9]%%chapternumber
}{}

\begin{document}
\maketitle
%\lehead{M. Beridze, N. Sharashenidze, L. Bakuradze \& Z. Pourtskhvanidze}

\section{Introduction}\label{3.8.ss.1}

Script, language, and cultural identity are often closely intertwined and play a crucial role in the demarcation and self-definition of linguistic communities, especially in isolated linguistic islands. This study examines how the \ili{Georgian} \isi{alphabet} has manifested itself as a central element in identity formation within the \ili{Georgian} language enclave in Iran (Fereydani \ili{Georgian}) and in the context of crypto-religiosity (Crypto-Orthodoxy, \citealt{pourtskhvanidze_beridze_bakuradze2023}) in the former \ili{Ottoman} area of Western Georgia. It addresses several key research questions. On the one hand, it examines how the \ili{Georgian} \isi{script} is used in these isolated communities to preserve and strengthen cultural and ethnic identity, and the role that isolation\is{Language!isolation} from the core \ili{Georgian} culture plays in this. It also analyses how this isolation\is{Language!isolation} affects the practice of writing. Another focus is on the influence of political and geographical isolation\is{Language!isolation}. The thesis sheds light on the historical development of the \ili{Georgian} language island and analyses the extent to which this isolation\is{Language!isolation} has contributed to the formation of an independent written and linguistic identity. The study examines how the Dedabruli\is{script|Dedabruli} \isi{script}, a variant of the \ili{Georgian} \isi{alphabet}, is used not only as a means of communication, but also as a symbol of cultural and religious identity, particularly in contexts where \ili{Georgian} communities were under Islamic influence. Finally, it discusses how writing and written traditions function as carriers of cultural memory and thus contribute to the formation of collective identity within the \ili{Georgian} language island in Iran. Specific written artefacts and their cultural significance will also be discussed. The study of these aspects will illustrate the interrelationship between writing, identity and linguistic isolation\is{Language!isolation}, highlighting the \ili{Georgian} \isi{alphabet} as a central identity-forming element in the Fereydani \ili{Georgian}.

\subsection{How does script shape identity?}\label{3.8.ss.1.1}

The identity construction of a group is based on a shared characteristic of the majority of this group, depending on the specific case \citep[2--3]{phinney2000}. Such a characteristic is contrastive to a shared characteristic of another group and thus fulfills the function of distinguishing or emphasizing the boundaries between the groups. We therefore refer to such shared characteristics as \textit{differentia specifica identitatis}. They motivate and enable the delimitation of a group from the inside and justify the corresponding perception from the outside. The \textit{differentia specifica identitatis} is a constitutive phenomenon of a group that regulates or balances the tension between isolation\is{Language!isolation} and integration with the help of these specific characteristics.

Which specific characteristic or combination of characteristics can be shared within a group to define identity? The list can be compiled with different rankings from such phenomena as language, religion, cultural memory, geography, customs and everyday traditions \citep[5]{segerstedt1942}. The dominance of a shared characteristic justifies the description of the group as a ``linguistic community'', ``religious community'', ``cultural community'', or ``geographical community'' \citep[4]{zabrocki1970}. 

The comparative analysis does not reveal a one-to-one correspondence between these ``communities'': a ``religious community'' may encompass multiple ``linguistic communities'', and, conversely, within a single linguistic community, diverse religious communities can emerge, establishing themselves as self-defining identity groups through shared religious affiliation. A similar constellation also applies to the ``scriptual community''.\footnote{%
We deliberately make no distinction here between \textit{writing} and \textit{script}\is{script}. The two phenomena merge in the given context and delimit the area in which the \textit{written image} (as a system and as an activity) is in the foreground. Consequently, this means equating handwriting (scriptography), book writing (typography), or even the use of letters for the purpose of counting.} This term is analogous to the term ``communicative community'' used by \citet{zabrocki1970} and goes back to the tradition of Rasmus Rask: ``[...] Language is based on human nature [...], writing is an accidental invention'' \citep{rask1826}. \citet[111--151]{glueck1987} writes ``[s]uch `arrangements' have taken very different historical forms and functions, and the respective writing communities were by no means necessarily congruent with the language communities that existed at the same time, as the examples of \ili{Latin} in the Middle Ages and early modern period or of \ili{English} in the world of the former Commonwealth show''. 

There is a type of multiple assignment between language and scriptual communities (\figref{ch3.8.chart1}).

\begin{figure}[t]
    \centering
    \includegraphics[width=.8\textwidth]{figures/fig9-1.pdf}
    \caption{Matrix of the multiple assignment between \textit{language community} and \textit{scriptual community}}
    \label{ch3.8.chart1}
\end{figure}

This captures the observation that one and the same language community can use different \isi{script} systems (e.g., \ili{Turkic} languages that utilize different \isi{script} systems in various \ili{Turkic}-speaking countries), as well as the fact that the same \isi{script} system can be used across different language communities (e.g., different languages from various language families represented in the \ili{Latin} \isi{script} system).

\largerpage[-2]
In the following article, a \textit{writing community}\is{writing community} is defined somewhat more broadly than the term ``writing'' is traditionally used. Here, the term refers to a community that defines itself as a community through the use of a unified and closed system of signs to reproduce linguistic content. The focus here is on the ``\ili{Georgian} language community'' and the ``\ili{Georgian} \isi{writing community}'' as one and the same projection area for both terms.\footnote{%
The \ili{Georgian} language belongs to the Kartvelian languages, which as a genetically closed language family (\ili{Georgian}, Zan and Svan) shows structural and typological similarities with other Caucasian languages. The general linguistic characteristics of the Kartvelian languages mainly concern the nominal and verbal morphology. \ili{Georgian}, for example, distinguishes between declination types of final-vowel-stem and final-consonant-stem. A special feature is the syntactic function of case. In the Kartvelian languages nominative, ergative and dative are regarded as subject cases. Ergative is the subject case in the transitive verbs and yet there are cases within the transitive verbs where the agent is not in the ergative but in the nominative. This typological peculiarity is called split-ergativity. Besides the complex nominal morphology, \ili{Georgian} has a very rich verbal morphology. Up to fifteen verb categories are listed. An important characteristic is the polypersonality. In the verb body, morphological markers mark not only the subject but also the direct and/or indirect object. An important role is played by the preverbs, whose function is not only the marking of the direction but also the formation of aspects and derivation. The grammatical tempus of the verb in \ili{Georgian} is a complex system of TAM categories
%\todo[inline]{Tense-Aspect-Mood ca}
and constructs three series. The sixteen dialects of \ili{Georgian} can initially be classified geographically into West \ili{Georgian} and East \ili{Georgian} dialects.} 

Traditionally, writing systems have been viewed in the context of identity formation in relation to religion \citep{haarmann1990}. Writing systems are seen primarily as a means of expressing religious identity. This includes the religiously determined choice of \isi{alphabet}, i.e. membership of a ``\isi{writing community}'' reflected membership of a ``religious community''. ``Jews wrote partly in \ili{Arabic}, but with \ili{Hebrew} letters. The first example is \textit{More Newuchim}, the major work of the Jewish scholar Moses ben Maimon, which was written in \ili{Arabic} but with \ili{Hebrew} letters'' \citep[524]{oppenheimer1971}.\footnote{%
Moreover, it is not only Jews but also Christians who underscore their religious identity through their choice of \isi{alphabet}: The Egyptians retained their language after Christianization in the first post-Christian centuries, but later adopted the \ili{Greek} \isi{alphabet} \citep[79]{metzger1966}. Atatürk had the \ili{Latin} \isi{script} introduced in 1928 as part of the \ili{Turkish} national reform, replacing the \ili{Arabic} \isi{script}. In this case, too, the writing system has the function of expressing one's identity -- only it is not a matter of religious identity, but of national identity.} However, religion, language, and \isi{script} are heterogeneous. The \ili{Arabic} \isi{script} is used not only to reproduce the \ili{Arabic} language, but also for other languages from other language families, such as \ili{Persian}, \ili{Kurdish} (in Iraq, Iran, and Syria), historically for \ili{Turkish} and \ili{Tatar}, Malay (Jawi), Indonesian, \ili{Pashto}, \ili{Urdu}, Uighur, Somali, Swahili, and Hausa.\footnote{%
Such uses of the \ili{Arabic} \isi{alphabet} for other languages are known as Ajami scripts.} In most instances, scripts are associated with the spread of Islam, particularly when non-indigenous writing systems are adopted for local languages. 

However, the case discussed in this article is unusual: here, the \ili{Georgian} \isi{script} is employed solely for the \ili{Georgian} language, rather than for writing other languages, as is more common in multilingual contexts influenced by religious or cultural exchange.\footnote{%
There are some minor exceptions, for example, in relation to the North Caucasian language Avar, which was written in the 14th century using the \ili{Georgian} \isi{alphabet} (\citealt[321--327]{chikobava1940}, \citealt[185--196]{gudava1954}). The inscription (tombstone) found in Khundzakhi in 1950 is a sentence in Avar (\textit{timala tacobeg} `[God] have mercy on the children') and is a mixture of two Mrglovani (round \isi{script}) and Nuskhuri (angular \isi{script}) variants of the \ili{Georgian} \isi{alphabet}.  The word nushuri is connected to the \ili{Georgian} word \georgian{ნუსხა} \textit{nuskh} `list, catalog', which in turn comes from the \ili{Arabic} word \textarab{نسخ} \textit{nasḫ} `copying'.} 
Thus, in this case, there is indeed a one-to-one relationship between the \textit{linguistic} and \textit{written} communities (Figure \ref{ch.3.8.chart2}). This fact makes it possible to see the \ili{Georgian} \isi{script} in the context of other identity-forming factors such as religion, cultural memory, architecture or traditional costume, although it should be noted that the order of emergence of identity-forming communities tends to favor the ``linguistic community'' before the others.\footnote{%
This assumption has historical roots in the \ili{Georgian} context. \textit{The Georgian Chronicles}\il{Georgian} (\georgian{ქართლის ცხოვრება} \textit{kartlis cxovreba}) (ch. 42.2)
%\todo[inline]{Add citation for Georgian Chronicles and Chapter 42.2} 
defines the Kingdom of Georgia (\georgian{ქართლის სამეფო} \textit{kartlis samepo}) through the linguistic practices that explain the idea of the identity of the kingdom in terms of religion:
\ea
\georgian{ქართლად ფრიადი ქუეყანაჲ აღირაცხების, რომელსაცა შინა ქართულითა ენითა ჟამი შეიწირვის და ლოცვაჲ ყოველი აღესრულების}\\
\textit{kartlad priadi kueq̇anay aġiracxebis, romelsaca šina kartulita enita žami šeic̣irvis da locvay q̇oveli aġesrulebis}\\
`Georgia is the country within which the service is celebrated in \ili{Georgian} and every prayer is requested'
\z}

The following article seeks to identify and describe those cases in which writing is appropriated not simply as a technological tool for graphically reproducing linguistic content, but rather as a medium for projecting beliefs and cultural or historical constructs -- analogous to expressions found in architecture, cuisine, or rites associated with marriage, baptism, and burial. 

\begin{figure}
    \centering
    \includegraphics[width=\textwidth]{figures/fig9-2.pdf}
    \caption{A one-to-one relationship between the \textit{linguistic} and \textit{writing} communities}
    \label{ch.3.8.chart2}
\end{figure}
%\todo[inline]{Check Y-axis label. Change "language community" to "linguistic community", or vice versa?}

A further distinction to be made at this point lies in the phenomenology of writing, which, unlike language, is the result of visual thinking. The connection between the phonetic image of language and the visuality of writing can be imagined on two levels: on the first primary level, the linguistic content is graphically visualized, while on the second level, based on this, writing represents the identity of the linguistic community and thus fulfills its secondary function (\textit{diferencia artificis itentitatis}). A further sub-level of visuality is represented by the written testimonies that represent the writing system itself: inscriptions, manuscripts, writing tools, writing practices, the scribes themselves, and related social networks. It should be noted that in the cases described below, the primary practice of using \ili{Georgian} \isi{script} for the purpose of reproducing linguistic content has taken a back seat, and the primary meaning has taken on a secondary function: the condition and anchoring of the identity of a linguistic community that represents a special form -- a language island.     

\largerpage
The following article deals with the cases of establishing (triggering) and representing identity through the \ili{Georgian} \isi{script} and related writing practices in the context of the \ili{Georgian} language islands.

\subsection{Georgian language islands and their origin: Brief overview}\label{3.8.ss.1.2}

The term \textit{language island} originates from \ili{German} linguistics, specifically dialectology, and was primarily employed in the early study of \isi{minority} languages and dialects \citep{ficker1869,kuhn1934}. Accordingly, language was initially the most important criterion for defining and describing the object of research. Of particular importance was the branch of research on language islands, which also included historical, cultural, and ethnic characteristics, and was thus located between linguistics, folklore studies/European ethnology, and history. A language island is defined as a segregated community of speakers that resists linguistic and cultural assimilation \citep{berend_mattheier1994}. 

Three \ili{Georgian} language islands, which exist outside present-day Georgia in Iran, Azerbaijan and Turkey, emerged as part of a historical continuum.\footnote{%
Georgia has experienced massive demographic change over the past approximately four hundred years. Historical processes covered by the term migration can be classified in various forms from prisoner-of-war migration to state-imposed ``demographic annexation''. The wars of conquest of the \ili{Persian} Shah Abbas the First in the first two decades of the 17th century, as well as the permanent invasions from the \ili{Ottoman} Empire turned large parts of the \ili{Georgian} Kingdom into deserted places. Shah Abbas continued his war policy with the means of demographic restructuring of Georgia: he resettled about 80 thousand families from mainly Kakheti (East Georgia) in the different locations of Iran and colonized the land thus freed by North Caucasian tribes of Muslim faith in the hope of achieving more loyalty to the \ili{Persian} Empire. In the course of this conflict, \ili{Georgian} language islands have emerged outside of \ili{Georgian} national territory: Fereydani \ili{Georgian} in the Islamic Republic of Iran, Ingilo \ili{Georgian} in Azerbaijan and the three variants spoken in Turkey: Taoish, Shavshish and Klarjian (on the fringes, the variants of the \ili{Georgian} Mujahir). The last of these is associated with the written form Dedabruli described in this article.} The armed conflict between the \ili{Ottoman} and \ili{Persian} Empires at the beginning and the Russian Empire in the later phase resulted in significant numbers of people from eastern \ili{Georgian} regions being deported to Iran (Fereydani \ili{Georgian}).\footnote{%
Depending on the source, the following places of origin can be identified: Kakhetia and Kiziq̇i (100,000 war slaves) \citep{rakhmani1957}; or Djavakhetia and Kartli (30,000 war slaves who were transferred to the villages north of Isfahan near \ili{Armenian} villages). According to \citet[16--46, 95--100, 102--104, 103--105]{davrizhetsi1669} (\armenian{Առաքել Դավրիժեցի} `Arakel of {Tabriz}'), the merchants from Tbilisi were also resettled in a suburb of Isfahan. From Tianetia (30,000 war slaves) and from Heretia (Saingilo) the Georgians were resettled to Mazandaran in Farahabad. After the conquest of the castle of Erevan, the shah's soldiers also reached Meskhetia and, according to Iskander Munsh, brought countless war slaves to Iran. According to Pierre \citet{oberling1963}, the Armenians and Georgians were resettled west of Isfahan almost simultaneously in the years 1603--5.} The resulting demographic annexation of East Georgia changed the linguistic landscape and formed a language island in the border region (Ingilo \ili{Georgian}). Russian interests in the South Caucasus and the wars with the \ili{Ottoman} Empire, which lasted until the beginning of the 20th century, separated the historical \ili{Georgian} cultural area in the southwest from the heartland and brought it within the newly founded \ili{Turkish} Republic. ``Our \ili{Georgian}'' (also ``Turkey \ili{Georgian}'') is an archipelago of different \ili{Georgian} islands with unclear borders in Turkey. All three language islands emerged from the \ili{Georgian} cultural area as part of the same political event that lasted over a century. All three islands emerged from different ethnicities and cultures but were in a similar Islamic-religious environment and were subject to different fates. 

\begin{figure}[htp]
     \centering
     \includegraphics[width=\textwidth]{figures/ch3.8.fig1.png}
     \caption{Mass forced resettlement under Shah Abbas the First. Illustration by the author.}
     \label{ch.3.8.fig.1}
 \end{figure}
%\todo[inline]{Check all instances of ``Illustration by the author'' and ``Photograph by the author''.}

There are two different possibilities or scenarios for the emergence of language islands:\footnote{%
This classification refers to the phenomena of the formation of \ili{Georgian} language islands, although language island phenomena outside the \ili{Georgian} language area (e.g. \ili{German} language islands) are also conceivable.} mass resettlement (see \figref{ch.3.8.fig.1}) and reduction of national territory. A distinction can be made between these two groups.

The other sub-forms of the emergence of language islands are:

\begin{enumerate}
    \item Forced resettlement during the war 
    \item Voluntary labor migration
    \item Forced labor migration (planned economy)
    \item Ecological catastrophes
\end{enumerate}

The reduction of national territory is rarely a peaceful process. It is usually the result of military confrontation with neighbors. An example of a mass forced resettlement is the Fereydani \ili{Georgian}. The second main scenario of the creation of language islands, the reduction of national territory, concerns Ingilo \ili{Georgian} -- a \ili{Georgian} language island in Azerbaijan.

\subsection{The case of Fereydani Georgian}\label{3.8.ss.1.3}

The linguistic isolation\is{Language!isolation} that occurred in the early 17th century led to the formation of Fereydani \ili{Georgian} \citep{bakuradze_beridze_pourtskhvanidze2020}. This is the Middle \ili{Georgian} phase in the development of the \ili{Georgian} language. Middle \ili{Georgian} begins in the late 12th century and ends in the 18th century. One hundred years after the forced resettlement of Eastern Georgians to Isfahan, the \ili{Georgian} language in Georgia has developed even further -- the new \ili{Georgian} language has begun to standardize its language through newly published newspapers, books and works in \ili{Georgian}, as well as the establishment of the \isi{National} Theatre, state schools and the \isi{Literacy} Society. The Georgians who were deported to Iran a century ago had no idea of the cultural development in their homeland. In the 19th century, Georgia made its way to Europe and all this was reflected in the ``new'' \ili{Georgian} language. The language of the Georgians in Fereydan differs from \ili{Georgian} by at least one stage of linguistic development. Shah Abbas I separated a part of the population from the main \ili{Georgian} nation, and consequently their language developed in a different way. The paths diverged for the next four hundred years, and only recently have they converged again \citep[494]{bakuradze_beridze_pourtskhvanidze2020}.

The language island, based on Fereydani \ili{Georgian}, can be defined as a geographically limited area where a language is used that differs from the same language or \isi{dialect} of the motherland by at least one historical stage of language development. 

The Fereydani-Georgians are at least bilingual (\ili{Persian}/Fereydani-\ili{Georgian}), and in villages where there is linguistic contact with South Luri, they speak three languages. In remote villages, there are many mixed families where parents teach their children both languages. When contact was re-established between Georgia and Fereydan, the Fereydani Georgians were very interested in learning the language and especially the \isi{alphabet}. \ili{Georgian} sources of the 17th and 18th centuries contain little information about the forced resettlement from Georgia to Iran. The first lasting contacts between the Georgians of Fereydan and \ili{Georgian} society began with the travelers from Fereydan. In the seventies of the 19th century, from time to time, \ili{Georgian} newspapers published news about \ili{Iranian} Georgians visiting the country. The first famous visitor from Fereydan was Atam Onikashvili, who published information about the life of \ili{Georgian} Fereydans in newspapers.  In 1922, Seipola Ioseliani visited Georgia. He was the first language assistant who helped to describe Fereydani \ili{Georgian}.

\section{Writing, identity and language islands}\label{3.8.ss.2}

In some cases, typographic signs can constitute meaning by integrating contextual elements of meaning into the overall text through the activation of culture-specific knowledge that need not be verbalized. 

\begin{figure}[t]
     \centering
     \includegraphics[width=.75\textwidth]{figures/ch3.8.fig2.png}
     \caption{Advertisement for the Austrian parliamentary elections, published in the \textit{Der Standard} newspaper on June 9, 1994 \citep[Figure 3]{schopp2008}}
     \label{ch.3.8.fig.2}
 \end{figure}

When the text element ``Others do it too'' in the advertisement (\figref{ch.3.8.fig.2}) for the 1994 Austrian parliamentary elections changes the font to Fraktur, the message is clear to every Austrian and \ili{German} with reference to the culturally specific relationship between broken fonts and Nazi ideology: ``The right-wing extremists are going to vote''. However, this shift in meaning is not understood in many parts of the world, where blackletter continues to bear the characteristic ``typically \ili{German}''. 

Especially in the context of language islands, the possibility of expressing an \isi{alphabet} is discussed not only from the perspective of the island language, but also in relation to the specific identity. The Transylvanian philologist Budai-Deleanu presents a comprehensive argument \citep{mitu2003}. His point is that the Cyrillic \isi{alphabet}, which is used in the surrounding Slavic languages, represents a foreign form for \ili{German}, which does not correspond at all to the internal structure of the language. There must be a close correspondence between \isi{alphabet} and language. Budai-Deleanu believes, in accordance with the linguistic theories of the time:
\begin{quote}
I could prove that the letters we call Church Slavonic are just as compatible with the inner essence and form of the Romanian language as the \ili{Turkish} ones. Because if we want to do violence to it, we can write it with \ili{Turkish} letters as well as with Slavic ones. But the question is: how will it deal with the other characteristics of the language? Because every language has its own. \citep[111]{budai-deleanu1970} 
\end{quote}
The author goes on to argue that this mismatch between language and \isi{alphabet} can lead to a serious problem of distortion of national identity: the Cyrillic \isi{alphabet} alienates one's Latinity, both from the eyes of foreigners and from one's own, and it alienates ``us'' from ``our'' own real condition. According to the author, only the reconstruction of the original orthographic garb can restore the ``transparency of identity''.

The activation of writing and literacy in the context of identity formation happens in different ways when it comes to language islands. One of the decisive factors is the spatial and institutional relationship to the heartland. For example, Volga Germans who emigrated from the Russian Empire to Pennsylvania in the 19th century received \ili{German} newspapers from Germany immediately after their resettlement \citep{mattheier1994}. 

Another factor is the social structures within the language island. At the time of the creation of the \ili{Georgian} language islands under consideration, it can be safely assumed that the idea, if not the active mastery and use of the \ili{Georgian} language, existed and was firmly established among the first generation of forced settlers. At present, as well as in the first travel accounts of the 19th century, there is no solid evidence of the active use of the \ili{Georgian} \isi{script} or whether it played a significant role at all.\footnote{%
In the travelogues of Agniashvili and Tchelidze from the 19th century, it is clear that there are questions about the possible practice of writing and written evidence. The assumption about the existence of hidden written artifacts was constantly verified and it was concluded that the Fereydani Georgians did not reveal them either out of fear or lack of trust towards the foreigners.}

\subsection{When script is intended to form identity: Fereydani Georgian}\label{3.8.ss.2.1}

What can be assumed, however, is the perception of the very different phonetic character of the \ili{Georgian} language, which partly exists as a language of memory, for example in some villages of the Fereydan. In the interviews with the language assistants in the village of Afus, the pharyngeal consonants of \ili{Georgian} were pointed out, and at the same time \ili{Persian} was set up as a contrastive background. It is likely that the perception of such differences in the active language communities creates a need for corresponding graphic representations. This would explain why in the last century there were at least three documented attempts by \ili{Georgian} Fereydans to create an \isi{alphabet} for \ili{Georgian} and introduce it into the life of the language island, without realizing that the \ili{Georgian} language has a writing tradition of almost two thousand years. Three hundred years of isolation\is{Language!isolation} from the mainland completely erased the idea of the writing tradition that was originally brought to the island. In an interview with one of the language assistants behind the last of the aforementioned attempts, it is reported that the goal of the new \isi{alphabet} creation was not primarily phonetic-phonological, but its integration into the everyday life of the Fereydani Georgians. The linguistic community created a special place for the phenomenon of writing in order to complement the \textit{differentia artificis identitatis}. 

At the end of the 19th century, after three hundred years of isolation\is{Language!isolation}, the reintegration of Fereydani \ili{Georgian} into the core \ili{Georgian} cultural area began. A geographically separate linguistic community began to become a scriptual community. The first travelers from Fereydan to the heartland indicate a clear mission: to travel to Georgia, learn to read and write in \ili{Georgian}, and then bring this knowledge back to Fereydan. The fact that there was no active or passive knowledge or memory of the writing tradition on the linguistic island itself at that time can be deduced indirectly from the records. In the transcripts of Agniashvili's dialogues \citep{agniashvili1896} with the Fereydani Georgians, they ask and talk about everything that concerns the heartland, but not a word about writing or literacy: ``[...] do you (in Georgia) have walnuts? Do you have plums? What do your chickens look like?''

Of course, there was no political or social opportunity to integrate the \ili{Georgian} \isi{script} into the life of the language island through the school system. Therefore, there was a deliberate search for and contact with Georgians who were sent from the Soviet Union near Fereydan (Najaf Abad, Isfahan) on various industrial projects (construction of a dam or a steel plant in Iran).\footnote{For help in interpreting the content, thanks are due to Prof. Giorgi Lobzhanidze and Dr. Tamar Lekveishvili.} 

\begin{figure}[t]
     \centering
     \includegraphics[width=\textwidth]{figures/ch3.8.fig4.png}
     \caption{Persian words written in Georgian script. /xomsdada/ ev. `to'/from' Xoms (/Xomsar); /guspand/(Pers) `sharp'; /xar/ (Pres) `donkey'; /nazomi/ (Fer Geo) `place'/`plot'; /rubendaz/ (Pers) `adiere (to it)'; /baq̇ie/ (Pers) `remainder.' After the numerals, letter /t/ probably stands for `toman.' Photograph by the author.}
     \label{ch.3.8.fig.4}
 \end{figure}
 %\todo[inline]{Invoice from 1918 with Persian words written in Georgian script. \georgian{ხომსდადამ} \textit{xomsdada} ‘to/from Xoms’ (\textit{Xomsar}); \georgian{გოსფანდ} \textit{guspand} (Pers \textarab{گوسفند}) ‘sheep’; \georgian{ხარ} \textit{xar} (Pers \textarab{خر}) ‘donkey’; \georgian{ნაზომი} \textit{nazomi} (Fereydani Geo) ‘place, plot’; \georgian{რუბენდაზ} \textit{rubendaz} (Pers) ‘adhered’; \georgian{ბაღიე} \textit{baq̇ie} (Per. \textarab{باقیه}) ‘remainder’. After the numerals, the letter \georgian{ტ} \textit{t} probably stands for ‘toman’ (a currency, Pers \textarab{تومان}). Photograph by the author. (Check spelling. Or \georgian{რუბენდაზ} \textit{rubendaz} (possibly Pers) ‘adhered (to it)’ or ‘recorded’)}

There are letters of the \ili{Georgian} \isi{alphabet} carved on the parquet from this period, along with words that have the corresponding phonemes as initial letters (for example, /ṗ/ in \georgian{პეპელა} \textit{ṗeṗela} `butterfly'). One artifact that came to light during fieldwork in Fereydan in 2018 is a 1918 invoice (\figref{ch.3.8.fig.4}), perhaps camouflaged by \ili{Georgian} \isi{script} (for whatever reason) and thus accessible only to certain readers.

The decisive step in the reintegration of Fereydani \ili{Georgian} into the core \ili{Georgian} cultural area is the formation of the written community from the heartland and the language island. 

In this process, two different phases can be distinguished: (1) personal contact with compatriots on business trips, which revives the cultural memory of the language island; (2) the state-organized distribution of \ili{Georgian} books in Fereydan on the part of the Soviet state. 

Another example from the field research of writing practices on the language island is the writing of one's own name (see \figref{ch.3.8.fig.5}).\footnote{For help in interpreting the content, thanks are due to Abolghasem Ahmadi.}

\begin{figure}[htp]
     \centering
     \includegraphics[width=\textwidth]{figures/ch3.8.fig5.png}
     \caption{Proper name \textit{karbalai alinaq̇i}. Photograph by the author.}
     \label{ch.3.8.fig.5}
\end{figure}
%\todo[inline]{Proper name in Georgian script: \georgian{ქარბალაი ალინაყი} \textit{karbalai alinaq̇i}, probably \textarab{كربلائي علي نقي} \textit{Karbalā’ī Ali Naqi} `Pilgrim Ali Naqi'. (check spelling)}

The \ili{Georgian} ethnicity of Joseph Stalin also served to revitalize identity in Fereydan and created a certain openness to Soviet policy, which aimed to bring Fereydani Georgians closer ideologically to the Soviet Union. The success of this policy required the active practice of literacy in the \ili{Georgian} language, i.e. the promotion of literacy using the \ili{Georgian} \isi{alphabet}. The archival material from this period documents extensive lists of party literature in the \ili{Georgian} language that were systematically sent to Fereydan. The motivation of the Fereydani Georgians for this policy was also due to the fact that they saw the policy of rapprochement as an opportunity to leave the island and return to their historical homeland, which in some cases actually happened. The letter allegedly written by Fereydani Georgians\footnote{%
Linguistically, the letter does not conform to the grammatical characteristics of Fereydani \ili{Georgian}. It can be assumed that the letter was written according to a model that did not originate with Fereydan.} to Stalin on April 30, 1944, which was found in the archives, is evidence of such sentiments (see \figref{ch.3.8.fig.6}).

\begin{figure}[htp]
     \centering
     \includegraphics[width=.9\textwidth]{figures/ch3.8.fig6.png}
     \caption{Archive of the Georgian Ministry of the Internal Affairs. Photograph by the author.}
     \label{ch.3.8.fig.6}
\end{figure}

In the houses and public places of Fereydan you can find a special tablecloth, which, in addition to the typical oriental ornaments, is bordered with a table (5 columns and 6 rows) containing the modern \ili{Georgian} \isi{alphabet} (\figref{ch.3.8.fig7}). 

\begin{figure}[htp]
     \centering
     \includegraphics[width=\textwidth]{figures/ch3.8.fig7.png}
     \caption{Letter to Stalin. Archives: Tablecloth with the alphabet cut out. Photograph by the author.}
     \label{ch.3.8.fig7}
\end{figure}
%\todo[inline]{Tablecloth \with Georgian alphabet and \georgian{ასოები ამოსაჭრელად} `alphabet to be cut out'.}

The last three boxes in the last row are connected and marked with the following sentence: ``Alphabet to be cut out''.
%\todo[inline{\georgian{ასოები ამოსაჭრელად} `Alphabet to be cut out' (check spelling)} 
When this sentence is applied to the tablecloth \isi{alphabet}, it seems illogical and even nonsensical, because if you follow the recommendation and cut out the individual letters from the tablecloth, the tablecloth will be destroyed (tablecloth cut up), and that cannot be the meaning of the sentence. Research at the manufacturer of this tablecloth in Fereydan revealed a story about the true event from the 1960s, which was embellished as a legend. The first Georgians from Fereydan who were officially invited to Soviet Georgia asked their hosts to take a \ili{Georgian} typewriter with them to Fereydan. Since all typewriters were registered in the Soviet Union and their ownership was excluded as private property, the \ili{Georgian} hosts could not possibly fulfill the wish of the guests from Fereydan. However, an employee of a printing house in Tbilisi procured a form of a table with the \ili{Georgian} \isi{alphabet}, which was used in Georgia for school beginners, so that they could receive duplicate sheets with the table, cut out the letters and thus acquire the \ili{Georgian} \isi{alphabet}. Now the form, which was apparently smuggled into Iran, was used to decorate a tablecloth.\footnote{%
In the story, the language assistant expressed the belief that this tablecloth means that the table is already set with spiritual food.}

Another story documented in the field research tells of a plastic bottle of mineral water from Georgia that was thrown into a garbage can by a field researcher at the Tehran airport. The collector, a Fereydani \ili{Georgian}, took the empty bottle from the trash can and carefully stored it. At home, the bottle was filled with ordinary table water and placed under the medicine cabinet -- there was a deep conviction that the name of the mineral water on the bottle label, inscribed with letters of the \ili{Georgian} \isi{alphabet}, would transform an ordinary table water into a remedy with healing powers through its divine effect. The liquid from the bottle was given to the sick and children in the hope of recovery or strengthening. 

In the context of language islands, the identity-forming function of the writing system can be supplemented by an additional aspect: the language of the language island, with its phonetic and grammatical characteristics, stands in maximum contrast to the language or languages of the surrounding area and constructs a projection surface for the foundation of identity. The writing system or the \ili{Georgian} \isi{alphabet} appears in the same light, which, like Fereydani \ili{Georgian}, stands in maximum contrast to the \ili{Arabic} writing system used in the surrounding \ili{Persian} language. The writing system takes on a sacred meaning for the \isi{maintenance} and preservation of identity. This meaning is expressed in everyday rituals and objects. The sacred meaning of the \ili{Georgian} \isi{alphabet} is even interpreted into existing facts that did not originally play a role in the context of the language's insular identity. An example of this is the bridge in Isfahan, which was built in the 17th century and is still an architectural feature of the city and the era. The bridge was built by the Fereydani \ili{Georgian} Alaverdi Khan Undiladze (Pers \textarab{اللّه وردی خان}) and is known as Siosepol (Pers \textarab{سیوسهپل}) `the bridge with thirty-three vaults'. In the Fereydani \ili{Georgian} narrative, this bridge is the symbol of the \ili{Georgian} \isi{alphabet}, as each individual arch of the bridge symbolizes one of the thirty-three letters of the \ili{Georgian} \isi{alphabet}. The Fereydani Georgians firmly believe that their famous compatriot built the bridge according to this principle. The reason why this is not true is the fact that the \ili{Georgian} \isi{alphabet} in the 17th century demonstrably contained 38 letters, not 33.

\subsection{Crypto-Orthodox crypto-alphabet: Dedabruli manuscripts}\label{3.8.ss.2.2}

%
Museums and archives in southwestern Georgia house numerous original 18th-century manuscripts known as ``dedabruli xeli'' \citep{khakhutashvili2017}. Translated as ``the old woman's hand'',
%\todo[inline]{Museums and archives in southwestern Georgia house numerous original 18th-century manuscripts written in a script known as \georgian{დედაბრული ხელი} \textit{dedabruli xeli} \citep{khakhutashvili2017}. Translated as `the old woman’s hand', (check spelling)} 
this term refers to the distinctive handwriting or calligraphy used by elderly women for clandestine communication among Orthodox \isi{minorities} in ``Osmalo, Georgia''.
%\todo[inline]{\georgian{ოსმალოს საქართველო} \textit{osmalos sakartvelo} `Ottoman Georgia'}
The earliest fragment of the \textit{Kobuletian-Adjarian Dedabruli Manuscripts} was discovered by D. Baqradze in 1873 (before the conflict between Russia and Turkey). In 1908, a letter from a Muhajir woman accompanied the Book of Amulets. In addition, private correspondence of Ruqie Romanadze (1900--1988), manuscripts with \ili{Arabic} incantations in the State Museum of Batumi, and the manuscript of Rukhset Surmanidze (died 1920) all date from the same year.
%

%
These manuscripts stand out because they were written by women. Linguists observe how mothers and grandmothers, who knew the \isi{script}, used it and passed on their knowledge to subsequent generations. Scholars of the \ili{Georgian} language try to draw parallels to this phenomenon in Western traditions \citep[11]{javakhishvili1949}.  
%

\begin{figure}[htp]
     \centering
     \includegraphics[width=\textwidth]{figures/ch3.8.fig8-cropped.png}
%     \includegraphics[width=\textwidth]{figures/ch3.8.fig10.png}
     \caption{The Dedabruli alphabet with the equivalents of the Georgian alphabet and transcription. Illustration by the author.}
     \label{ch.3.8.fig.8}
\end{figure}
%\todo[inline]{Use figures/ch3.8.fig10.png instead. The Dedabruli alphabet with the equivalents of the Georgian alphabet and transliteration. Illustration by the author.}

They refer to the so-called ``nun alphabets'' (/monazvnuri/ from (Geo) /monazoni/ `nun') and note that the ``dedabruli xeli'' is also referred to as the ``monazvnuri'' \isi{script}.
%\todo[inline]{They refer to the so-called ``nun alphabets'' (\georgian{მონაზვნური} \textit{monazvnuri} `nun-related' from \georgian{მონაზონი} \textit{monazoni} `nun') and note that the \textit{dedabruli xeli} is also referred to as the \textit{monazvnuri} \isi{script}. (check spelling)} 
What is fascinating about this story is the intersection between the activities of the nuns and the secret communication facilitated by an invented \isi{alphabet}.

The \isi{script} found in these manuscripts is a fusion of two \ili{Georgian} alphabets, deliberately distorted graphically to render them unreadable without a deep understanding of specific graphic conventions (see \figref{ch.3.8.fig.8}).

Some letters are graphically similar to the original letters and have recognizable references to each other (examples: /z/, /w/). Some letters have been adopted unchanged from the modern \ili{Georgian} \isi{alphabet} (examples: /č̣/, /č/, /w/). In /p/ and /ǯ/, for example, the handwritten versions have been changed in such a way that the graphic reference is no longer recognizable. Especially noteworthy is /ž/. 
%\todo[inline]{Some letters are graphically similar to the original letters and have recognizable references to each other (examples: \georgian{ზ} \textit{z}, \georgian{ვ} \textit{w}). Some letters have been adopted unchanged from the modern \ili{Georgian} \isi{alphabet} (examples: \georgian{ჭ} \textit{č̣}, \georgian{ჩ} \textit{č}, \georgian{ვ} \textit{w}). In \georgian{ფ} \textit{p} and \georgian{ჯ} \textit{ǯ}, for example, the handwritten versions have been changed in such a way that the graphic reference is no longer recognizable. Especially noteworthy is \georgian{ჟ} \textit{ž}. (Why repeat \georgian{ვ} \textit{w}? Especially noteworthy is \georgian{ო} \textit{o}. \georgian{ჟ} \textit{ž} is missing.}
Another difficulty is that some letters have a different phonetic value, known only to the correspondents, in order to increase the degree of ciphering. The \ili{Georgian} version of the language game\footnote{%
Known in \ili{German} as \textit{Löffelsprache}.} is also used in spelling, where an additional vowel or syllable is consistently inserted between the syllables of a word. Word division was not observed, and punctuation marks were dropped in all manuscripts.

Lexically, there are some so-called ``abracadabra'' words, where the knowledge of their meaning is no longer used in treatises and which are written after remembering the sound: /prayer against freezing: \textsuperscript{(?)}\textit{ila} \textsuperscript{(?)}\textit{ilma} \textsuperscript{(?)}\textit{napal} \textsuperscript{(?)}\textit{mana gadebian ǯinikma}/.
%\todo[inline]{sound: prayer against freezing, /[ʔ]ila [ʔ]ilma [ʔ]napal [ʔ]mana gadebian ǯinikma/.}

Khakhunashvili \citep[553]{khakhunashvili_tsetskhladze2019} claims that this writing tradition emerged in Orthodox monasteries in the late Middle Ages and later took on a secret communication function during the \ili{Ottoman} Empire. The term ``old woman's hand'' probably developed later to refer to a form of calligraphy that was considered crude because of its inconsistent graphic representation. This interpretation is supported by the words of one scribe: ``I scribbled it in the manner of old women''. Scholars use the term /dedabruli xeli/ ``old women's hand'' to describe these manuscripts, suggesting a tradition that predates the \ili{Ottoman} conquest of Christian areas in Georgia. 
%\todo[inline]{The term \textit{dedabruli xeli} `old woman's hand' probably developed later to refer to a form of calligraphy that was considered crude because of its inconsistent graphic representation. This interpretation is supported by the words of one scribe: ``I scribbled it in the manner of old women''. Scholars use the term \textit{dedabruli xeli} to describe these manuscripts, suggesting a tradition that predates the \ili{Ottoman} conquest of Christian areas in Georgia.}
Nevertheless, this style of writing took on a new role in shaping and preserving identity during this historical period. Written by women and passed down through generations, these manuscripts served different purposes. Some were in the form of long strips that could be folded and concealed in amulets or talismans, containing prayers and gospel texts. Others contained private information about relocated family members. 

Overall, the aim of the \textit{dedabruli xeli} manuscripts was to strengthen the cultural identity of the Orthodox \isi{minority} under \ili{Ottoman} rule and to preserve the memory of the scriptural and literary traditions of \ili{Georgian} Orthodoxy. It is worth noting that some of these manuscripts are incomplete or completely illegible (see \figref{ch.3.8.fig.9}).

\begin{figure}
     \centering
     \includegraphics[height=.9\textheight]{figures/ch3.8.fig9.png}
     \caption{Examples of scrolls with Dedabruli script for the amulets or talismans from the Kobuleti Manuscript Museum (Georgia). The long strips of paper hidden in the talismans list the churches and holy places bearing the name of St. George, the patron saint of Georgia. Photograph by the author.}
     \label{ch.3.8.fig.9}
 \end{figure}

The Dedabruli \isi{alphabet} can be seen in the context of the preservation of religious identity. The forcibly Islamised population preserves its ethnic identity as a linguistic community, but sees its religious identity threatened and practices a so-called Crypto-Orthodox Christianity. They use a writing system derived from the traditional \ili{Georgian} \isi{alphabet}, which is only accessible to certain people (usually women). The Dedabruli \isi{alphabet} is of great cultural and linguistic importance for the \ili{Georgian} language island in Turkey. The areas influenced by Turkey were larger in the 18th and 19th centuries than the present-day area of \georgian{ჩვენებურები} \textit{Chveneburebi} `Turkey Georgians' lit. `of us'. 
%\todo{check spelling}
The area of the spread of the Dedabruli writing was included. This community, which lived under Islamic rule, used the Dedabruli \isi{script} as an instrument to preserve their religious (Crypto-Orthodox), ethnic and linguistic identity. The \isi{script} offered them the opportunity to maintain a sense of cultural continuity despite the overwhelming influence of Islam. The \textit{Chveneburebi} relied on the Dedabruli \isi{alphabet} not only as a practical writing system, but also as a symbol of their \ili{Georgian} heritage. The \isi{alphabet} became a sacred representation of their cultural memory and identity, and strengthened their distinctiveness from \ili{Turkish} culture.

\section{Conclusion}\label{3.8.ss.3}

In the \ili{Georgian} linguistic community, writing plays a much larger role than simply representing the language -- it becomes a central feature of cultural identity. But what does it mean when a writing system not only serves to communicate, but also contributes to the formation of identity? Particularly in isolated linguistic communities, such as the \ili{Georgian} communities in Iran, Turkey and Azerbaijan, the \ili{Georgian} \isi{script} is used as a ``secondary function'' to strengthen cultural cohesion. Here it plays a crucial role in differentiating from foreign cultures and assisting the community in its cultural self-assertion. But why is the ``secondary function'' so important for the survival of cultural identity?

The primary function of a writing system is to represent language visually -- it is the means by which phonetic elements are recorded graphically. But in these \ili{Georgian} language islands, writing has taken on a completely different role: it serves as a symbol, an anchor to preserve and strengthen cultural memory and ethnic consciousness in a foreign cultural environment. But why is this necessary? Why resort to a system of writing in an alien environment that is alien to the inhabitants of that environment and apparently has no practical function? This ``secondary function'' describes the ability of writing to do much more than convey information. It embodies the cultural heritage of the community and provides Fereydani Georgians with a link to their roots. The \isi{script} becomes a conscious marker of identity, distinct from the dominant \ili{Persian} culture and its \ili{Arabic} \isi{script}. By cultivating and preserving the \ili{Georgian} \isi{script}, even when there is no practical reason to do so, Fereydani Georgians express a strong sense of their cultural origins. The \isi{script} becomes a symbol of \isi{belonging} to the \ili{Georgian} nation and maintains a distinction from their surroundings. 

This symbolic use of the \isi{script} goes beyond mere writing. It is integrated into everyday life, into rituals and even into everyday objects such as talismans or religious texts. \ili{Georgian} writing thus becomes a vehicle for meanings that represent not only the word but an entire cultural consciousness. These rituals show that this is not a writing that is only read -- it is experienced, celebrated, revered. The writing itself acquires a quasi-sacred status and is charged with meanings that are deeply embedded in cultural identity. It creates a cultural space that exists independently of practical communication. By incorporating these signs of the \ili{Georgian} \isi{script} into everyday objects and symbolic artifacts, the Fereydani Georgians make the writing a living symbol of their identity. The Fereydani-Georgians see the \ili{Georgian} \isi{script} as a cultural bond that connects them to their homeland and distinguishes them from the \ili{Persian}-speaking environment. Thus, in these isolated language communities, writing becomes a unifying element that strengthens the community's identity and builds a cultural bridge. But why do the Fereydani-Georgians continue to cling to a writing system that is not necessary for communication in their environment? It allows them to preserve their identity; for the language island, \ili{Georgian} writing becomes a means of resisting assimilation and the loss of their own identity.

The dominance of the secondary function described here justifies the use of the term \textit{differentia specifica identitatis} in relation to the \isi{script}. In writing, certain aesthetic, cultural and religious ideas are interpreted or acted out. On this level, writing behaves similarly to all other surfaces onto which cultural, ethnic and religious identity is projected. In the case of a language island, a particular intensity of the secondary function of language can be assumed. The common language connects the language island, and the common writing system strengthens this connection.

\section*{Abbreviations}
\begin{tabularx}{.5\textwidth}{@{}lQ@{}}
TAM & Tense-Aspect-Mood \\
Geo & Georgian \\
\end{tabularx}%
\begin{tabularx}{.5\textwidth}{@{}lQ@{}}
Pers & Persian \\
& \\
\end{tabularx}

\section*{Acknowledgements}
This work was supported by the Shota Rustaveli \isi{National} Science Foundation of Georgia (SRNSFG) [grant number FR-21-336].

\sloppy
\printbibliography[heading=subbibliography,notkeyword=this]
\end{document}
